\section{Explaining the models}
\label{explain}

\subsection{Logistic-regression coefficients}

\textbf{Global view.} Table~\ref{tab:terms} lists the largest positive and negative coefficients
after refitting the selected structured-LR configuration on seed-0 training rows. Positive terms
increase the model's fraud log-odds when other feature values are held fixed. Correlated TF-IDF terms
share weight, so these are associations within one fitted model, not independent or causal effects.

\begin{table}[tb]
\centering
\small
\setlength{\tabcolsep}{4pt}
\begin{tabular}{lrlr}
\toprule
\multicolumn{2}{c}{\textbf{Towards fraud}} & \multicolumn{2}{c}{\textbf{Towards real}} \\
\textbf{Feature} & \textbf{Coef.} & \textbf{Feature} & \textbf{Coef.} \\
\midrule
\tabinput{includes/tab_terms}
\bottomrule
\end{tabular}
\caption{Largest structured-LR coefficients on seed-0 training rows, ten in each direction.}
\label{tab:terms}
\end{table}

\textbf{Interpretation.} Money-related words (\emph{money}, \emph{income}, \emph{earn}) coexist with
placeholder/link phrases such as \emph{link url} and \emph{apply using}. Industry and function terms
also have large weights. This mixture is consistent with both scam language and dataset-specific
poster/category patterns; the table alone cannot separate those explanations.

\subsection{Decision-tree importances}

\textbf{Global view.} Table~\ref{tab:tree-terms} reports impurity-based feature importances for the
selected seed-0 tree. Company-profile presence dominates with importance 0.389, followed by the
\emph{High School or Equivalent} education category at 0.116 and company-logo presence at 0.037. The strong
profile effect agrees descriptively with the class-conditional difference in
Figure~\ref{fig:eda-flags}.

\begin{table}[tb]
\centering
\small
\begin{tabular}{lr}
\toprule
\textbf{Feature} & \textbf{Importance} \\
\midrule
\tabinput{includes/tab_tree_terms}
\bottomrule
\end{tabular}
\caption{Ten largest impurity-based feature importances for the selected seed-0 decision tree.}
\label{tab:tree-terms}
\end{table}

\textbf{Boundary.} Impurity importance describes splits in this training-fitted tree and can favour
features that offer useful training partitions. It is not a held-out permutation-importance
estimate. Both classical explanation tables belong to one inspected partition and do not establish
that the same features transfer to another split, source or time period.

\subsection{Errors and neural evidence}

\textbf{Classical error inspection.} Notebook \texttt{02} prints the highest-scoring validation false
alarms and lowest-scoring validation misses for structured LR. Examples are selected on validation
rather than test errors, separating the displayed cases from test labels.

\textbf{Current neural evidence.} The neural analyses reported here are behavioral comparisons of
text-only/fused inputs and top-two/full adaptation budgets, plus an input-token coverage audit. No
neural occlusion, integrated-gradients, attention-attribution or counterfactual explanation analysis
has been run. Fine-tuning does compute gradients for optimization, but those gradients are not
reported as token explanations.

\textbf{Interpretation limit.} The fine-tuned CLS state has no direct one-to-one word
interpretation. The classical coefficients and tree importances are therefore the only global
feature-attribution artifacts presented, while the paired neural results indicate predictive
performance, not which linguistic evidence the neural models use.
